\section[K. Aeonglass]{\myboxxi{      \BF{\TT{K}}} Aeonglass}

\textit{Time limit}   2 seconds

\textit{Memory limit} 1 Gigabyte

\textit{Input} \qquad standard input

\textit{Output} \qquad standard output


\begin{quote}
	\textit{任季节一再轮转，我却始终没能习惯。}
\end{quote}

\subsection{Problem Statement}

In the Euclidean plane, you are given a polygon $P=P_0P_1\cdots P_{n-1}$ whose vertices are listed in counterclockwise order.

An angle $\alpha$ ($0 \le \alpha < \pi$) is a \BF{stable direction} if and only if there exists a $c \in \R$ such that, for every point $X$ in polygon $P$ (including its boundary), the perpendicular segment from $X$ to the line
\[
	\cos \alpha\,x + \sin \alpha\,y + c = 0
\]
lies entirely in $P$ (including its boundary). Here $\alpha$ is a direction of the line's normal vector.


Find the number of \BF{stable directions}. This number can be shown to be finite.

\subsection{Input}

The first line contains an integer $T$ ($1 \le T \le 100$), the number of test cases.

The first line of each test case contains an integer $n$ ($3 \le n \le 10^5$), the number of polygon vertices.

Each of the next $n$ lines contains two integers $x_i,y_i$ ($-10^9 \le x_i,y_i \le 10^9$), the coordinates of the $i$-th vertex, listed in counterclockwise order.

Within a test case, all vertices are distinct. Adjacent edges intersect only at their shared endpoint, and nonadjacent edges do not intersect.

The sum of $n$ over all test cases does not exceed $10^5$.

\subsection{Output}

For each test case, output one integer: the number of \BF{stable directions}.

\subsection{Sample 1}

\begin{minipage}{\linewidth}
	\begin{minipage}[t]{0.49\linewidth}
		\begin{lstlisting}[title=\texttt{standard input}]
3
4
0 0
1 0
2 1
1 1
4
0 0
2 0
2 2
0 2
6
0 0
4 0
4 1
1 1
3 4
0 4

\end{lstlisting}
	\end{minipage}
	\begin{minipage}[t]{0.49\linewidth}
		\begin{lstlisting}[title=\texttt{standard output}]
3
4
1
\end{lstlisting}
	\end{minipage}
\end{minipage}

\newpage
\subsection{Sample 1 Explanation}

% \paragraph{样例示意图}

\begin{figure}[h]
	\centering

	% ============================================================
	% Test 1
	% ============================================================
	\begin{minipage}[t]{0.31\linewidth}
		\centering
		\begin{tikzpicture}[
			scale=1.25,
			>=Stealth,
			every node/.style={font=\scriptsize}
		]
			% axes
			\draw[->] (-0.45,0) -- (2.55,0)
				node[below] {$x$};
			\draw[->] (0,-0.45) -- (0,1.75)
				node[left] {$y$};

			% ticks
			\foreach \x in {1,2}
				\draw (\x,0.04) -- (\x,-0.04)
					node[below=2pt] {\x};
			\draw (0.04,1) -- (-0.04,1)
				node[left=2pt] {1};

			% polygon
			\fill[Line11!50]
				(0,0) -- (1,0) -- (2,1) -- (1,1) -- cycle;
			\draw[Line11!50!black, very thick]
				(0,0) -- (1,0) -- (2,1) -- (1,1) -- cycle;

			% x = 1
			\draw[black!70, very thick, dashed]
				(1,-0.25) -- (1,1.45);
			\node[black!70, right] at (1.03,1.35)
				{$L_1$};

			% x + y = 3/2
			\draw[black!70, very thick, dashed]
				(0.05,1.45) -- (1.75,-0.25);
			\node[black!70, above right]
				at (0.10,1.38)
				{$L_2$};

			% x - 2y = 0
			\draw[black!70, very thick, dashed]
				(-0.25,-0.125) -- (2.35,1.175);
			\node[black!70, above]
				at (1.85,0.93)
				{$L_3$};

			% vertices
			\fill (0,0) circle (1.1pt);
			\fill (1,0) circle (1.1pt);
			\fill (2,1) circle (1.1pt);
			\fill (1,1) circle (1.1pt);
		\end{tikzpicture}

		\vspace{1mm}

	\t{\scriptsize Illustration for test case 1}
	\end{minipage}
	\hfill
	% ============================================================
	% Test 2
	% ============================================================
	\begin{minipage}[t]{0.31\linewidth}
		\centering
		\begin{tikzpicture}[
			scale=1.05,
			>=Stealth,
			every node/.style={font=\scriptsize}
		]
			% axes
			\draw[->] (-0.45,0) -- (2.65,0)
				node[below] {$x$};
			\draw[->] (0,-0.45) -- (0,2.65)
				node[left] {$y$};

			% ticks
			\foreach \x in {1,2}
				\draw (\x,0.04) -- (\x,-0.04)
					node[below=2pt] {\x};
			\foreach \y in {1,2}
				\draw (0.04,\y) -- (-0.04,\y)
					node[left=2pt] {\y};

			% polygon
			\fill[Line11!50]
				(0,0) -- (2,0) -- (2,2) -- (0,2) -- cycle;
			\draw[Line11!50!black, very thick]
				(0,0) -- (2,0) -- (2,2) -- (0,2) -- cycle;

			% x = 1
			\draw[black!70, very thick, dashed]
				(1,-0.25) -- (1,2.4);
			\node[black!70, right] at (1.04,2.25)
				{$L_1$};

			% y = 1
			\draw[black!70, very thick, dashed]
				(-0.25,1) -- (2.4,1);
			\node[black!70, below right]
				at (2.05,1.03)
				{$L_2$};

			% x + y = 2
			\draw[black!70, very thick, dashed]
				(-0.15,2.15) -- (2.15,-0.15);
			\node[black!70, above right]
				at (0.10,1.92)
				{$L_3$};

			% x - y = 0
			\draw[black!70, very thick, dashed]
				(-0.15,-0.15) -- (2.25,2.25);
			\node[black!70, below right]
				at (2.05,2.05)
				{$L_4$};

			% vertices
			\fill (0,0) circle (1.1pt);
			\fill (2,0) circle (1.1pt);
			\fill (2,2) circle (1.1pt);
			\fill (0,2) circle (1.1pt);
		\end{tikzpicture}

		\vspace{1mm}

	\t{\scriptsize Illustration for test case 2}
	\end{minipage}
	\hfill
	% ============================================================
	% Test 3
	% ============================================================
	\begin{minipage}[t]{0.34\linewidth}
		\centering
		\begin{tikzpicture}[
			scale=0.72,
			>=Stealth,
			every node/.style={font=\scriptsize}
		]
			% axes
			\draw[->] (-0.55,0) -- (4.75,0)
				node[below] {$x$};
			\draw[->] (0,-0.55) -- (0,4.75)
				node[left] {$y$};

			% ticks
			\foreach \x in {1,2,3,4}
				\draw (\x,0.06) -- (\x,-0.06)
					node[below=2pt] {\x};
			\foreach \y in {1,2,3,4}
				\draw (0.06,\y) -- (-0.06,\y)
					node[left=2pt] {\y};

			% polygon
			\fill[Line11!50]
				(0,0) -- (4,0) -- (4,1) --
				(1,1) -- (3,4) -- (0,4) -- cycle;
			\draw[Line11!50!black, very thick]
				(0,0) -- (4,0) -- (4,1) --
				(1,1) -- (3,4) -- (0,4) -- cycle;

			% unique valid line
			\draw[black!70, very thick, dashed]
				(0.5,-0.3) -- (0.5,4.4);
			\node[black!70, right] at (0.56,4.35)
				{$L$};

			% vertices
			\fill (0,0) circle (1.2pt);
			\fill (4,0) circle (1.2pt);
			\fill (4,1) circle (1.2pt);
			\fill (1,1) circle (1.2pt);
			\fill (3,4) circle (1.2pt);
			\fill (0,4) circle (1.2pt);

			% highlight reflex vertex
			% \node[above right] at (1,1)
			% 	{$(1,1)$};
		\end{tikzpicture}

		\vspace{1mm}

	\t{\scriptsize Illustration for test case 3}
	\end{minipage}

	% \vspace{4em}
	% 三组样例及其合法直线示意图。虚线表示可以选取的直线。
\end{figure}

The lines and \BF{stable directions} for test case 1 are

\[
\begin{aligned}
L_1:\quad &x-1=0,
&\alpha_1&=0,\\
L_2:\quad &\frac{\sqrt{2}}{2}x+\frac{\sqrt{2}}{2}y-\frac{3\sqrt{2}}{4}=0,
&\alpha_2&=\frac{\pi}{4},\\
L_3:\quad &-\frac{1}{\sqrt{5}}x+\frac{2}{\sqrt{5}}y=0,
&\alpha_3&=\pi-\arctan 2.
\end{aligned}
\]



The lines and \BF{stable directions} for test case 2 are

\[
\begin{aligned}
L_1:\quad &x-1=0,
&\alpha_1&=0,\\
L_2:\quad &y-1=0,
&\alpha_2&=\frac{\pi}{2},\\
L_3:\quad &\frac{\sqrt{2}}{2}x+\frac{\sqrt{2}}{2}y-\sqrt{2}=0,
&\alpha_3&=\frac{\pi}{4},\\
L_4:\quad &-\frac{\sqrt{2}}{2}x+\frac{\sqrt{2}}{2}y=0,
&\alpha_4&=\frac{3\pi}{4}.
\end{aligned}
\]

The line and \BF{stable direction} for test case 3 are

\[
L:\quad x-\frac{1}{2}=0,
\qquad
\alpha_1=0.
\]

The corresponding line is not uniquely determined.
